% page 218 of the facsimile (image p-217 of the scan)
% block 165.1 x 250.6 mm ; left margin 8.0 mm
\pgc{-12.700mm}{0.000mm}{
\l{\cel{3}210}
\l{}
\l{\sou{heleboro}.\cel{2}(\sou{bot.)}\cel{2}Planto herbacea de la familio \textquotedbl{}ranunkula-}
\l{\cel{3}cei\textquotedbl{}, di qua la stipi esas rametoza, la flori verdatra,}
\l{\cel{3}e la radiko uzata medicinale por la morbi nervala. -L.}
\l{\cel{3}helleborus. - EFISL.}
\l{}
\l{helio.\cel{2}(kemio).\cel{2}Gaso di atmosfero, deskovrita sur Suno (per}
\l{\cel{3}la analizo spektrala) e pose sur Tero (en uranio-erco).}
\l{\cel{3}-\cel{1}\sou{Simbolo kemiala}\cel{1}: He.\cel{2}- DEFIRS.}
\l{}
\l{\sou{helianto}.\cel{2}(\sou{bot.)}\cel{2}Genero de planti de la familio \textquotedbl{}kompozaji\textquotedbl{},}
\l{\cel{3}a qua apartenas la sunfloro e la topinamburo. -L.Helianthus.}
\l{\cel{3}- EFIS.}
\l{}
\l{\sou{heliantemo}. (\sou{bot.})\cel{2}Genero de planti ek la familio \textquotedbl{}cistacei\textquotedbl{}.}
\l{\cel{3}- L.\cel{1}\sou{helianthemum}.}
\l{}
\l{\sou{helico.}\cel{2}I. (\sou{geom.}) Lineo du-kurva, trasita filetatre cirkum}
\l{\cel{3}cilindro. - II. (\sou{mekaniko)}.\cel{2}Aparato helicatra qua, movata}
\l{\cel{3}da vaporo o motoro, propulsas la navi, la aeroplani, e c.}
\l{\cel{3}EFIS.}
\l{}
\l{\sou{helicoido}. (\sou{geom.}) Surfaco produktita da rekto perpendikla ye}
\l{\cel{3}la axo vertikala di cilintro rekta, e turnanta alonge la}
\l{\cel{3}helico trasita sur la cilindro.}
\l{}
\l{.heliko. (\sou{zool.}) Molusko gasteropoda, reptera, di qua la konko}
\l{\cel{3}esas helicatra. - L.\cel{1}\sou{cochlea}.}
\l{}
\l{\sou{helikoptero}. Aviac-aparato qua su elevas preske vertikale en at-}
\l{\cel{3}mosfero (e decensas de ibe) per surfaci rotacanta quin movas}
\l{\cel{3}motoro. - EFIRS.}
\l{}
\l{\sou{heliocentrala}. (\sou{astron.}) Qua relatas la centro di Suno. - DEF.}
\l{}
\l{\sou{heliofilo}. Qua prizas la lumo-radii.}
\l{}
\l{\sou{heliofoba}. Qua patologie timas la lumo-radii.}
\l{}
\l{\sou{heliograbar}. (\sou{trans.}) Reproduktar imaji fotografura per imprimo}
\l{\cel{3}de sika graburo.}
\l{}
\l{\sou{heliografar}. (\sou{trans}.) Aplikar la procedi per qui reproduktesas,}
\l{\cel{3}sur paperfolii fotografala, kalquo-desegnuri di qui la linei}
\l{\cel{3}esas nigra sur fundo blanka.}
\l{}
\l{\sou{heliometro}. (\sou{astron.})\cel{2}Aparato uzata por mezurar la diametro}
\l{\cel{3}semblanta di Suno e di la planeti. - DEFIRS.}
\l{}
\l{\sou{helioskopo}. (\sou{astron.}) Vitro-peco koloroza per qua onu povas}
\l{\cel{3}regardar Suno. - DEFIRS.}
\l{}
\l{\sou{heli}ostato. (\sou{tekn.}) Aparato optikala qua mantenas en konstanta}
\l{\cel{3}direciono sunoradio enduktita en kamero obskura. - DEFIRS.}
\l{}
\l{heliotropo.(bot.)\cel{2}Planto de la familio \textquotedbl{}boraginei\textquotedbl{}, di qua un}
\l{\cel{3}speco havas flori qui odoras dolcege. L. Heliotropium europeum}
}
